\par
\SourceLeaf{242}
I. Si $M(a,b,c)$ est R. Q. numeri $m$, hic determinantem formae $(a,b,c)$ metietur. Si enim $(g,h)$ est valor expressionis $\surd M(a,b,c)$ (mod.\ $m$), sive
\[
gg\equiv aM,\quad gh\equiv bM,\quad hh\equiv cM\ (\mathrm{mod.}\ m)
\]
erit $bbMM-acMM\equiv0$, sive $(bb-ac)MM$ per $m$ divisibilis. Quoniam autem $M$ ad $m$ primus esse supponitur, etiam $bb-ac$ per $m$ divisibilis erit.

II. Si $M(a,b,c)$ est R. Q. ipsius $m$, atque $m$ aut numerus primus aut potestas numeri primi, puta $=p^\mu$: character particularis formae $(a,b,c)$ respectu numeri $p$ erit vel $Rp$ vel $Np$, prout $M$ est residuum vel non-residuum ipsius $p$. Hoc statim inde sequitur, quod tum $aM$ tum $cM$ est residuum ipsius $m$ sive ipsius $p$, atque ad minimum unus numerorum $a$, $c$ per $p$ non divisibilis (art. 230).

Simili modo, si (manentibus reliquis) $m=4$, erit vel $1,4$ vel $3,4$ character part. formae $(a,b,c)$, prout $M\equiv1$ vel $\equiv3$; nec non si $m=8$ vel altior potestas numeri $2$. erit $1,8;\ 3,8;\ 5,8;\ 7,8$ char. part. formae $(a,b,c)$, prout $M\equiv1;\ 3;\ 5;\ 7\ (\mathrm{mod.}\ 8)$ resp.

III. Vice versa si $m$ est numerus primus aut numeri primi imparis potestas $=p^\mu$, determinantem $bb-ac$ metiens, atque $M$ vel residuum vel non-residuum ipsius $p$, prout character formae $(a,b,c)$ respectu ipsius $p$ est $Rp$ vel $Np$ resp., erit $M(a,b,c)$ resid. quadr. ipsius $m$. Quando enim $a$ per $p$ non est divisibilis, $aM$ erit res. ipsius $p$ adeoque etiam ipsius $m$; si itaque $g$ est valor expr. $\surd aM$ (mod.\ $m$), $h$ valor expr. $\frac{bg}{a}$ (mod.\ $m$), erit $gg\equiv aM;\ ah\equiv bg$, adeoque
\[
agh\equiv bgg\equiv abM\quad\text{et}\quad gh\equiv bM
\]
denique
\[
ahh\equiv bgh\equiv bbM\equiv bbM-(bb-ac)M\equiv acM
\]
adeoque $hh\equiv cM$, \textit{i. e.}\ $(g,h)$ valor expr. $\surd M(a,b,c)$. Quando vero $a$ per $m$ est divisibilis, certo $c$ non erit; unde facile perspicitur, eadem resultare, si pro $h$ assumatur valor expr. $\surd cM$ (mod.\ $m$). pro $g$ valor expr. $\frac{bh}{c}$ (mod.\ $m$).

Simili modo demonstratur, si $m$ fuerit $=4$ ipsumque $bb-ac$ metiatur, numerusque $M$ accipiatur vel $\equiv1$ vel $\equiv3$, prout $1,4$ vel $3,4$ fuerit char. part.
\SourceLeaf{243}
formae $(a,b,c)$: fore $M(a,b,c)$ res. qu. ipsius $m$. Nec non, si $m$ fuerit $=8$ vel altior potestas ipsius $2$, per quam $bb-ac$ divisibilis sit, atque $M$ accipiatur $\equiv1;\ 3;\ 5;\ 7\ (\mathrm{mod.}\ 8)$, prout character part. formae $(a,b,c)$ respectu numeri $8$ postulet: $M(a,b,c)$ fore res. qu. ipsius $m$.

IV. Si determinans formae $(a,b,c)$ est $=D$, atque $M(a,b,c)$ res. qu. ipsius $D$, omnes characteres particulares formae $(a,b,c)$ tum respectu singulorum divisorum primorum imparium ipsius $D$, tum respectu numeri $4$ vel numeri $8$ (si ipsum $D$ metiuntur) ex numero $M$ statim cognosci possunt. Ita \textit{e. g.}\ quum $3(20,10,27)$ sit resid. qu. ipsius $440$, scilicet $(150,9)$ valor expr. $\surd3(20,10,27)$ sec. mod.\ $440$, atque $3N5$, $3R11$: characteres formae $(20,10,27)$ sunt $3,8;\ N5;\ R11$. Soli characteres particulares respectu numerorum $4$ et $8$, quoties determinantem non metiuntur, nexum necessarium cum numero $M$ non habent

V. Vice versa, si numerus $M$ ad $D$ primus omnes characteres particulares formae $(a,b,c)$ in se complectitur (exceptis characteribus respectu numerorum $4,8$, quando ipsum $D$ non metiuntur): erit $M(a,b,c)$ res. qu. ipsius $D$. Nam ex III patet, si $D$ sub formam $\pm A^\alpha B^\beta D^\gamma\ldots$ redigatur, ita ut $A$, $B$, $C$ etc. sint numeri primi diversi, fore $M(a,b,c)$ resid. qu. singulorum $A^\alpha$, $B^\beta$, $C^\gamma$ etc. Si igitur valor expr. $\surd M(a,b,c)$ secundum mod.\ $A^\alpha$, est $(\mathfrak{A},\mathfrak{A}')$; secundum mod.\ $B^\beta$, $(\mathfrak{B},\mathfrak{B}')$; sec. mod.\ $C^\gamma$, $(\mathfrak{C},\mathfrak{C}')$ etc. numerique $g$, $h$ ita determinantur ut sit $g\equiv\mathfrak{A},\mathfrak{B},\mathfrak{C}$ etc.; $h\equiv\mathfrak{A}',\mathfrak{B}',\mathfrak{C}'$ etc. secundum modulos $A^\alpha$, $B^\beta$, $C^\gamma$ etc. resp. (art. 32): facile perspicietur, fore $gg\equiv aM$, $gh\equiv bM$, $hh\equiv cM$ secundum omnes modulos $A^\alpha$, $B^\beta$, $C^\gamma$ etc. adeoque etiam secundum modulum $D$ qui illorum est productum.

VI. Propter has rationes numeri tales ut $M$ vocabuntur \textit{numeri characteristici} formae $(a,b,c)$, poteruntque per V plures huiusmodi numeri nullo negotio inveniri, simulac omnes characteres particulares huius formae sunt eruti; simplicissimi autem tentando plerumque evolvuntur facillime. Manifestum est, si $M$ sit numerus characteristicus formae primitivae datae determinantis $D$, omnes numeros, ipsi $M$ secundum mod.\ $D$ congruos, fore numeros characteristicos eiusdem formae; formas in eadem classe, sive etiam in classibus diversis ex eodem genere, contentas eosdem numeros characteristicos habere, quamobrem quivis numerus
\SourceLeaf{244}
characteristicus formae datae etiam toti classi et generi tribui potest; denique $1$ semper esse numerum characteristicum formae classis et generis principalis, sive quamlibet formam e genere principali esse residuum determinantis sui.

VII. Si $(g,h)$ est valor expr. $\surd M(a,b,c)$ (mod.\ $m$), atque $g'\equiv g$, $h'\equiv h$ (mod.\ $m$): erit etiam $(g',h')$ valor eiusdem expressionis. Tales valores pro \textit{aequivalentibus} haberi possunt; contra si $(g,h)$, $(g',h')$ sunt valores eiusdem expr. $\surd M(a,b,c)$, neque tamen simul $g'\equiv g$, $h'\equiv h$ (mod.\ $m$), \textit{diversi} sunt censendi. Manifesto quoties $(g,h)$ est valor talis expressionis, etiam $(-g,-h)$ erit, facileque demonstratur, hos valores semper esse diversos nisi $m=2$. Aeque facile demonstratur, expressionem $\surd M(a,b,c)$ (mod.\ $m$) plures valores diversos quam duos tales (oppositos) habere non posse, quando $m$ sit aut numerus primus impar aut numeri primi imparis potestas aut $=4$; quando vero $m$ sit $=8$ aut altior potestas numeri $2$, quatuor omnino dari. Hinc facile deducitur per VI, si determinans $D$ formae $(a,b,c)$ sit $=\pm2^\mu A^\alpha B^\beta\ldots$, designantibus $A$, $B$ etc. numeros primos impares diversos quorum multitudo $=n$, atque $M$ numerus characteristicus illius formae: dari omnino vel $2^n$ vel $2^{n+1}$ vel $2^{n+2}$ valores diversos expr. $\surd M(a,b,c)$ (mod.\ $D$), prout $\mu$ vel $<2$ vel $=2$ vel $>2$. Ita \textit{e. g.}\ habentur sedecim valores expr. $\surd7(12,6,-17)$ (mod.\ $240$), puta $(\pm18,\mp11)$, $(\pm18,\pm29)$, $(\pm18,\mp91)$, $(\pm18,\pm109)$, $(\pm78,\pm19)$, $(\pm78,\pm59)$, $(\pm78,\mp61)$, $(\pm78,\mp101)$. Demonstrationem ampliorem quum ad sequentia non sit adeo necessaria, brevitatis gratia non apponimus.

VIII. Denique observamus, si duarum formarum aequivalentium $(a,b,c)$, $(a',b',c')$ determinans sit $D$, numerus characteristicus $M$, priorque transeat in posteriorem per substitutionem $\alpha,\beta,\gamma,\delta$: ex quovis valore expr. $\surd M(a,b,c)$ ut $(g,h)$ sequi valorem expr. $\surd M(a',b',c')$, puta $(\alpha g+\gamma h,\beta g+\delta h)$. Demonstrationem quisque nullo negotio eruere poterit.

\Needspace{12\baselineskip}
\begin{center}
\textit{De compositione formarum.}
\end{center}

\Needspace{8\baselineskip}
\Article{234}
Postquam haec de formis in classes genera et ordines distribuendis praemisimus, proprietatesque generales quae ex his distinctionibus statim defluunt explicavimus, ad aliud argumentum gravissimum transimus a nemine hucusque
\SourceLeaf{245}
attactum, de formarum \textit{compositione}. In cuius disquisitionis limine, ne posthac demonstrationum seriem continuam interrumpere oporteat, statim intercalamus

\textsc{Lemma.} \textit{Habentur quatuor series numerorum integrorum}
\[
\begin{gathered}a,a',a''\ldots a^n;\qquad b,b',b''\ldots b^n;\\c,c',c''\ldots c^n;\qquad d,d',d''\ldots d^n\end{gathered}
\]
\textit{ex aeque multis (puta $n+1$) terminis constantes, atque ita comparatae, ut}
\[
cd'-dc',\quad cd''-dc''\ \text{\textit{etc.}},\quad c'd''-d'c''\ \text{\textit{etc. etc.}}
\]
\textit{respective sint}
\[
=k(ab'-ba'),\quad k(ab''-ba'')\ \text{\textit{etc.}},\quad k(a'b''-b'a'')\ \text{\textit{etc. etc.}}
\]
\textit{sive generaliter}
\[
c^\lambda d^\mu-d^\lambda c^\mu=k(a^\lambda b^\mu-b^\lambda a^\mu)
\]
\textit{denotante $k$ numerum integrum datum; $\lambda$, $\mu$ integros quoscunque inaequales inter $0$ et $n$ incl. quorum maior $\mu$}\HistoricalNote{*)}{Considerando $a$ tamquam $a^0$, $b$ tamquam $b^0$ etc. --- Ceterum manifesto eadem aequatio valebit quoque quando $\lambda=\mu$ aut $\lambda>\mu$.}\textit{; praeterea omnes $a^\lambda b^\mu-b^\lambda a^\mu$ divisorem communem non habent. Tunc inveniri possunt quatuor numeri integri $\alpha,\beta,\gamma,\delta$ tales, ut sit}
\[
\begin{aligned}
\alpha a+\beta b&=c, &\alpha a'+\beta b'&=c', &\alpha a''+\beta b''&=c''\ \text{\textit{etc.}}\\
\gamma a+\delta b&=d, &\gamma a'+\delta b'&=d', &\gamma a''+\delta b''&=d''\ \text{\textit{etc.}}
\end{aligned}
\]
\textit{sive generaliter}
\[
\alpha a^\nu+\beta b^\nu=c^\nu,\quad\gamma a^\nu+\delta b^\nu=a^\nu
\]
\textit{quo facto erit}
\[
\alpha\delta-\beta\gamma=k
\]

Quum per hyp. numeri $ab'-ba'$, $ab''-ba''$ etc. $a'b''-b'a''$ etc. (quorum multitudo erit $=\frac12(n+1)n$) divisorem communem non habeant, inveniri poterunt totidem alii numeri integri, per quos illis resp. multiplicatis productorum summa fiat $=1$ (art. 40). Designentur hi multiplicatores per $(0,1)$, $(0,2)$ etc. $(1,2)$ etc., sive generaliter multiplicator ipsius $a^\lambda b^\mu-b^\lambda a^\mu$ per $(\lambda,\mu)$, ita ut sit
\[
\Sigma(\lambda,\mu)(a^\lambda b^\mu-b^\lambda a^\mu)=1
\]
(Per literam $\Sigma$ denotamus aggregatum omnium valorum expressionis, cui praefixa
\SourceLeaf{246}
est, qui oriuntur tribuendo ipsis $\lambda$, $\mu$ omnes valores inaequales inter $0$ et $n$, ita ut sit $\mu>\lambda$). Quo facto si statuitur
\[
\begin{aligned}
\Sigma(\lambda,\mu)(c^\lambda b^\mu-b^\lambda c^\mu)&=\alpha, &\Sigma(\lambda,\mu)(a^\lambda c^\mu-c^\lambda a^\mu)&=\beta\\
\Sigma(\lambda,\mu)(d^\lambda b^\mu-b^\lambda d^\mu)&=\gamma, &\Sigma(\lambda,\mu)(a^\lambda d^\mu-d^\lambda a^\mu)&=\delta
\end{aligned}
\]
hi $\alpha,\beta,\gamma.\delta$ proprietatibus praescriptis erunt praediti.

\textit{Dem.} I. Denotante $\nu$ numerum quemcunque integrum inter $0$ et $n$, erit
\[
\begin{aligned}
\alpha a^\nu+\beta b^\nu
&=\Sigma(\lambda,\mu)(c^\lambda b^\mu a^\nu-b^\lambda c^\mu a^\nu+a^\lambda c^\mu b^\nu-c^\lambda a^\mu b^\nu)\\
&=\frac1k\Sigma(\lambda,\mu)(c^\lambda d^\mu c^\nu-d^\lambda c^\mu c^\nu)\\
&=\frac1k c^\nu\Sigma(\lambda,\mu)(c^\lambda d^\mu-d^\lambda c^\mu)\\
&=c^\nu\Sigma(\lambda,\mu)(a^\lambda b^\mu-b^\lambda a^\mu)=c^\nu
\end{aligned}
\]
Et per calculum similem eruitur
\[
\gamma a^\nu+\delta b^\nu=d^\nu.\qquad\text{\textit{Q. E. P.}}
\]

II. Quoniam igitur
\[
c^\lambda=\alpha a^\lambda+\beta b^\lambda,\quad c^\mu=\alpha a^\mu+\beta b^\mu
\]
fit
\[
c^\lambda b^\mu-b^\lambda c^\mu=\alpha(a^\lambda b^\mu-b^\lambda a^\mu)
\]
similique modo
\[
\begin{aligned}
a^\lambda c^\mu-c^\lambda a^\mu&=\beta(a^\lambda b^\mu-b^\lambda a^\mu)\\
d^\lambda b^\mu-b^\lambda d^\mu&=\gamma(a^\lambda b^\mu-b^\lambda a^\mu)\\
a^\lambda d^\mu-d^\lambda a^\mu&=\delta(a^\lambda b^\mu-b^\lambda a^\mu)
\end{aligned}
\]
ex quibus formulis valores ipsorum $\alpha,\beta,\gamma,\delta$ multo facilius erui possunt, si modo $\lambda$, $\mu$ ita accipiuntur, ut $a^\lambda b^\mu-b^\lambda a^\mu$ non sit $=0$, quod certo fieri poterit, quia omnes $a^\lambda b^\mu-b^\lambda a^\mu$ per hyp. divisorem communem non habent, adeoque \textit{omnes} $=0$ esse nequeunt. --- Ex iisdem aequationibus deducitur, multiplicando primam per quartam, secundam per tertiam et subtrahendo,
\[
(\alpha\delta-\beta\gamma)(a^\lambda b^\mu-b^\lambda a^\mu)^2
=(a^\lambda b^\mu-b^\lambda a^\mu)(c^\lambda d^\mu-d^\lambda c^\mu)
=k(a^\lambda b^\mu-b^\lambda a^\mu)^2
\]
unde necessario
\[
\alpha\delta-\beta\gamma=k.\qquad\text{\textit{Q. E. S.}}
\]

\Needspace{8\baselineskip}
\SourceLeaf{247}
\Article{235}
Si forma
\[
AXX+2BXY+CYY\ldots F
\]
transit in productum e duabus formis
\[
axx+2bxy+cyy\ldots f,\quad\text{et}\quad a'x'x'+2b'x'y'+c'y'y'\ldots f'
\]
per substitutionem talem
\[
\begin{aligned}
X&=pxx'+p'xy'+p''yx'+p'''yy'\\
Y&=qxx'+q'xy'+q''yx'+q'''yy'
\end{aligned}
\]
(quod brevitatis causa in sequentibus semper ita exprimemus: Si $F$ transit in $ff'$ per substitutionem $p,p',p'',p''';\ q,q',q'',q'''$\HistoricalNote{*)}{In hac igitur designatione ad ordinem tum coefficientium $p,p'$ etc. tum formarum $f,f'$ probe respicere oportet. Facile autem perspicietur, si ordo formarum $f,f'$ convertatur ut prior fiat posterior coefficientes $p',q'$ cum his $p'',q''$ commutandos esse, reliquos suo quemlibet loco manere.}), dicemus simpliciter, formam $F$ \textit{transformabilem} esse in $ff'$; si insuper haec transformatio ita est comparata, ut sex numeri
\[
\begin{gathered}pq'-qp',\quad pq''-qp'',\quad pq'''-qp''',\\p'q''-q'p'',\quad p'q'''-q'p''',\quad p''q'''-q''p'''\end{gathered}
\]
divisorem communem non habeant: formam $F$ e formis $f,f'$ \textit{compositam} vocabimus.

Inchoabimus hanc disquisitionem a suppositione generalissima, formam $F$ in $ff'$ transire per substitutionem $p,p',p'',p''';\ q,q',q'',q'''$ et quae inde sequantur evolvemus. Manifesto huic suppositioni ex asse aequivalebunt sequentes novem aequationes (\textit{i. e.}\ simulac hae aequationes locum habent, $F$ per substitutionem dictam transibit in $ff'$, et vice versa):
\begin{align*}
App+2Bpq+Cqq&=aa'\quad\ldotp\,\ldotp\,\ldotp\tag*{[1]}\\
Ap'p'+2Bp'q'+Cq'q'&=ac'\quad\ldotp\,\ldotp\,\ldotp\tag*{[2]}\\
Ap''p''+2Bp''q''+Cq''q''&=ca'\quad\ldotp\,\ldotp\,\ldotp\tag*{[3]}\\
Ap'''p'''+2Bp'''q'''+Cq'''q'''&=cc'\quad\ldotp\,\ldotp\,\ldotp\tag*{[4]}\\
App'+B(pq'+qp')+Cqq'&=ab'\quad\ldotp\,\ldotp\,\ldotp\tag*{[5]}\\
App''+B(pq''+qp'')+Cqq''&=ba'\quad\ldotp\,\ldotp\,\ldotp\tag*{[6]}\\
Ap'p'''+B(p'q'''+q'p''')+Cq'q'''&=bc'\quad\ldotp\,\ldotp\,\ldotp\tag*{[7]}\\
Ap''p'''+B(p''q'''+q''p''')+Cq''q'''&=cb'\quad\ldotp\,\ldotp\,\ldotp\tag*{[8]}\\
\begin{aligned}A(pp'''+p'p'')&+B(pq'''+qp'''+p'q''+q'p'')\\&+C(qq'''+q'q'')\end{aligned}&=2bb'\quad\ldotp\,\ldotp\,\ldotp\tag*{[9]}
\end{align*}

\SourceLeaf{248}
Sint determinantes formarum $F,f,f'$ resp. $D,d,d'$; divisores communes maximi numerorum $A,2B,C;\ a,2b,c;\ a',2b',c'$ resp. $M,m,m'$ (quos omnes positive acceptos supponimus). Porro determinentur sex numeri integri $\mathfrak{A},\mathfrak{B},\mathfrak{C},\mathfrak{A}',\mathfrak{B}',\mathfrak{C}'$ ita ut sit
\[
\mathfrak{A}a+2\mathfrak{B}b+\mathfrak{C}c=m,\qquad
\mathfrak{A}'a'+2\mathfrak{B}'b'+\mathfrak{C}'c'=m'
\]
Denique designentur numeri
\[
\begin{gathered}pq'-qp',\quad pq''-qp'',\quad pq'''-qp''',\\p'q''-q'p'',\quad p'q'''-q'p''',\quad p''q'''-q''p'''\end{gathered}
\]
resp. per $P,Q,R,S,T,U$, sitque ipsorum divisor communis maximus positive acceptus $=k$. --- Iam ponendo
\begin{align*}
App'''+B(pq'''+qp''')+Cqq'''&=bb'+\Delta\quad\ldotp\,\ldotp\,\ldotp\,\ldotp\tag*{[10]}
\end{align*}
fit ex aequ. 9
\begin{align*}
Ap'p''+B(p'q''+q'p'')+Cq'q''&=bb'-\Delta\quad\ldotp\,\ldotp\,\ldotp\,\ldotp\tag*{[11]}
\end{align*}

Ex his undecim aequationibus $1\ldots11$, sequentes novas evolvimus\HistoricalNote{*)}{Origo harum aequationum haec est: 12 ex $5.5-1.2$; 13 ex $5.9-1.7-2.6$; 14 ex $10.11-6.7$; 15 ex $5.8+5.8+10.10+11.11-1.4-2.3-6.7-6.7$; 16 ex $8.9-3.7-4.6$; 17 ex $8.8-3.4$. Deductio sex reliquarum eodem modo adornatur, si modo aequationes. $2,5,7$ cum aequationibus $3,6,8$ resp. commutantur, et reliquae $1,4,9,10,11$ eodem loco deinceps retinentur, puta 18 ex $6.6-1.3$ etc.}:
\begin{align*}
DPP&=d'aa\quad\ldotp\,\ldotp\,\ldotp\,\ldotp\,\ldotp\,\ldotp\,\ldotp\,\ldotp\,\ldotp\,\ldotp\,\ldotp\,\ldotp\,\ldotp\tag*{[12]}\\
DP(R-S)&=2d'ab\quad\ldotp\,\ldotp\,\ldotp\,\ldotp\,\ldotp\,\ldotp\,\ldotp\,\ldotp\,\ldotp\,\ldotp\,\ldotp\,\ldotp\tag*{[13]}\\
DPU&=d'ac-(\Delta\Delta-dd')\quad\ldotp\,\ldotp\,\ldotp\,\ldotp\,\ldotp\,\ldotp\,\ldotp\tag*{[14]}\\
D(R-S)^2&=4d'bb+2(\Delta\Delta-dd')\quad\ldotp\,\ldotp\,\ldotp\,\ldotp\,\ldotp\tag*{[15]}\\
D(R-S)U&=2d'bc\quad\ldotp\,\ldotp\,\ldotp\,\ldotp\,\ldotp\,\ldotp\,\ldotp\,\ldotp\,\ldotp\,\ldotp\,\ldotp\,\ldotp\tag*{[16]}\\
DUU&=d'cc\quad\ldotp\,\ldotp\,\ldotp\,\ldotp\,\ldotp\,\ldotp\,\ldotp\,\ldotp\,\ldotp\,\ldotp\,\ldotp\,\ldotp\,\ldotp\tag*{[17]}\\
DQQ&=da'a'\quad\ldotp\,\ldotp\,\ldotp\,\ldotp\,\ldotp\,\ldotp\,\ldotp\,\ldotp\,\ldotp\,\ldotp\,\ldotp\,\ldotp\tag*{[18]}\\
DQ(R+S)&=2da'b'\quad\ldotp\,\ldotp\,\ldotp\,\ldotp\,\ldotp\,\ldotp\,\ldotp\,\ldotp\,\ldotp\,\ldotp\,\ldotp\tag*{[19]}\\
DQT&=da'c'-(\Delta\Delta-dd')\quad\ldotp\,\ldotp\,\ldotp\,\ldotp\,\ldotp\,\ldotp\tag*{[20]}\\
D(R+S)^2&=4db'b'+2(\Delta\Delta-dd')\quad\ldotp\,\ldotp\,\ldotp\,\ldotp\,\ldotp\tag*{[21]}\\
D(R+S)T&=2db'c'\quad\ldotp\,\ldotp\,\ldotp\,\ldotp\,\ldotp\,\ldotp\,\ldotp\,\ldotp\,\ldotp\,\ldotp\,\ldotp\,\ldotp\tag*{[22]}\\
DTT&=dc'c'\quad\ldotp\,\ldotp\,\ldotp\,\ldotp\,\ldotp\,\ldotp\,\ldotp\,\ldotp\,\ldotp\,\ldotp\,\ldotp\,\ldotp\,\ldotp\tag*{[23]}
\end{align*}

Hinc rursus deducuntur hae duae:
\SourceLeaf{249}
\[
\begin{aligned}
0&=2d'aa(\Delta\Delta-dd')\\
0&=(\Delta\Delta-dd')^2-2d'ac(\Delta\Delta-dd')
\end{aligned}
\]
scilicet prior ex $12.15-13.13$, posterior ex $14.14-12.17$; unde facile perspicitur, necessario esse $\Delta\Delta-dd'=0$, sive sit $a=0$, sive non sit $=0$\HistoricalNote{*)}{Haec derivatio aequationis $\Delta\Delta=dd'$ ad institutum praesens sufficit; alioquin analysin elegantiorem sed hic nimis prolixam tradere possemus, directe deducendo ex aequationibus $1\ldots11$ hanc $0=(\Delta\Delta-dd')^2$.}. Supponemus itaque, in aequatt. $14,15,20,21$ ad dextram deleri $\Delta\Delta-dd'$.

Iam statuendo
\[
\begin{aligned}
\mathfrak{A}P+\mathfrak{B}(R-S)+\mathfrak{C}U&=mn'\\
\mathfrak{A}'Q+\mathfrak{B}'(R+S)+\mathfrak{C}'T&=m'n
\end{aligned}
\]
(ubi $n,n'$ etiam fractiones evadere posse probe notandum, etsi $mn',m'n$ necessario sint integri): facile ex aequatt. $12\ldots17$ deducitur
\[
Dmmn'n'=d'(\mathfrak{A}a+2\mathfrak{B}b+\mathfrak{C}c)^2=d'mm
\]
similiterque ex aequ. $18\ldots23$
\[
Dm'm'nn=d(\mathfrak{A}'a'+2\mathfrak{B}'b'+\mathfrak{C}'c')^2=dm'm'
\]

Erit igitur $d=Dnn,\ d'=Dn'n'$, unde nanciscimur \textsc{conclusionem primam}: \textit{Determinantes formarum $F,f,f'$ necessario inter se habent rationem quadratorum;} et \textsc{secundam}: \textit{$D$ semper metitur numeros $dm'm',d'mm$.} Patet itaque, $D,d,d'$ eadem signa habere, nullamque formam in productum $ff'$ transformabilem esse posse, cuius determinans maior sit quam divisor communis maximus numerorum $dm'm',d'mm$.

Multiplicentur aequationes $12,13,14$ resp. per $\mathfrak{A},\mathfrak{B},\mathfrak{C}$; similiterque per eosdem numeros aequatt. $13,15,16$, et $14,16,17$; addantur terna producta, dividaturque summa per $Dmn'$, scripto pro $d'$, $Dn'n'$. Tunc prodit
\[
P=an',\quad R-S=2bn',\quad U=cn'
\]

Simili modo multiplicatis aequationibus $18,19,20$ nec non $19,21,22$ et $20,22,23$ resp. per $\mathfrak{A}',\mathfrak{B}',\mathfrak{C}'$, obtinetur
\[
Q=a'n,\quad R+S=2b'n,\quad T=c'n
\]

\SourceLeaf{250}
Hinc habetur \textsc{conclusio tertia}: \textit{Numeri $a,2b,c$ proportionales sunt numeris $P,R-S,U$, positaque illorum ratione ad hos ut 1 ad $n'$, erit $n'$ radix quadrata ex $\frac{d'}{D}$; similiterque numeri $a',2b',c'$ ad $Q,R+S,T$ eandem rationem habent, quae si ponitur esse ut 1 ad $n$, erit $n$ radix quadrata ex $\frac{d}{D}$.}

Ceterum quantitates $n,n'$ radices vel positivae vel negativae ex $\frac{d}{D},\frac{d'}{D}$ esse possunt, unde distinctionem petimus, quae primo aspectu sterilis videbitur, sed cuius usus in sequentibus sufficienter apparebit. Scilicet dicemus, in transformatione formae $F$ in $ff'$ formam $f$ accipi \textit{directe} quando $n$ est positiva, \textit{inverse} quando $n$ negativa; similiterque $f'$ accipi directe vel inverse, prout $n'$ positiva vel negativa. Accedente autem conditione ut $k$ sit $=1$, forma $F$ \textit{vel} ex utraque forma $f,f'$ directe composita, \textit{vel} ex utraque inverse \textit{vel} ex $f$ directe et ex $f'$ inverse, \textit{vel} ex $f$ inverse et ex $f'$ directe dicetur, prout vel $n,n'$ ambae sunt positivae, vel ambae negativae, vel prior positiva posterior negativa, vel prior negativa posterior positiva. Ceterum quisque facile intelliget, has relationes ab ordine quo formae $f,f'$ collocantur (vid. annot. prim. ad art. praes.) non pendere.

Porro observamus, divisorem maximum communem numerorum $P,\allowbreak Q,\allowbreak R,\allowbreak S,\allowbreak T,\allowbreak U$ puta $k$ metiri numeros $mn',m'n$ (uti ex valoribus supra stabilitis manifestum est) adeoque quadratum $kk$ ipsos $mmn'n',m'm'nn$, atque $Dkk$ ipsos $d'mm,dm'm'$. Sed et vice versa quivis divisor communis ipsorum $mn',m'n$ metietur ipsum $k$. Sit enim $e$ talis divisor, qui manifesto etiam numeros $an',2bn',cn',a'n,2b'n,c'n$ metietur, \textit{i. e.} numeros $P,R-S,U,Q,R+S,T$ et proin etiam ipsos $2R$ et $2S$. Iam si $\frac{2R}{e}$ esset numerus impar, etiam $\frac{2S}{e}$ impar esse deberet (quoniam summa et differentia sunt pares) adeoque etiam productum impar. Hoc autem productum fit $=\frac{4}{ee}(b'b'nn-bbn'n')=\frac{4}{ee}(d'nn+a'c'nn-dn'n'-acn'n')=\frac{4}{ee}(a'c'nn-acn'n')$ adeoque par, quia $e$ ipsos $a'n,c'n,an',cn'$ metitur. Quare $\frac{2R}{e}$ necessario erit par, et proin $R$ nec non $S$ per $e$ divisibilis. Quoniam igitur $e$ omnes sex $P,Q,R,S,T,U$ metitur, metietur etiam ipsorum divisorem communem maximum $k$. \textit{Q. E. D.} --- Hinc concluditur, $k$ esse divisorem communem maximum numerorum $mn',m'n$; unde facile perspicietur, \textit{$Dkk$ fore divisorem communem maximum numerorum $dm'm',d'mm$.} Quae est \textsc{conclusio quarta}. Patet itaque, quoties $F$ ex $f$ et $f'$ composita sit, $D$ fore divisorem communem maximum, numerorum $dm'n',d'mm$, et vice versa; quae proprietas etiam tamquam definitio formae compositae adoptari potuisset. Forma igitur composita e
\SourceLeaf{251}
formis $f,f'$ determinantem maximum possibilem inter omnes formas in productum $ff'$ transformabiles habet.

Antequam ulterius progredi possimus, ante omnia valorem ipsius $\Delta$ accuratius definire oportet, quem quidem ostendimus esse $=\surd dd'=\surd DDnnn'n'$, sed cuius \textit{signum} hinc nondum determinatur. Ad hunc finem ex aequatt. fundamentalibus 1---11 eruimus $DPQ=\Delta aa'$ (quae aequ. obtinetur ex $5.6-1.11$), adeoque $Daa'nn'=\Delta aa'$, unde, nisi aliquis numerorum $a,a'$ est $=0$, fit $\Delta=Dnn'$. Sed prorsus simili modo ex aequatt. fundd. octo aliae deduci possunt, in quibus ad laevam $Dnn'$ ad dextram $\Delta$ multiplicati habeantur per $2ab',ac',2ba',4bb',2bc',ca',2cb',cc'$\HistoricalNote{*)}{Analysin quam lectores facile detegere poterunt brevitatis caussa supprimere oportet.}, unde facile concluditur propterea quod neque omnes $a,2b,c$, neque omnes $a',2b',c'$ possunt esse $=0$, in omnibus casibus fieri $\Delta=Dnn'$, adeoque $\Delta$ idem signum habere ut $D,d,d'$ vel oppositum, prout $n,n'$ eadem signa habeant vel diversa.

Porro observamus, numeros $aa',2ab',ac',2ba',4bb',2bc',ca',2cb',cc',2bb'+2\Delta,2bb'-2\Delta$ omnes per $mm'$ divisibiles esse. De novem prioribus hoc per se manifestum est, de duobus reliquis autem simili modo demonstrari potest ut antea ostendimus $R$ et $S$ per $e$ divisibiles esse. Scilicet patet, $4bb'+4\Delta$ et $4bb'-4\Delta$ per $mm'$ divisibiles esse (quoniam $4\Delta=\surd16dd'$ atque $4d$ per $mm$, $4d'$ per $m'm'$ divisibilis, adeoque $16dd'$ per $mmm'm'$ et $4\Delta$ per $mm'$) et differentiam quotientium parem; productum ex quotientibus facile demonstratur esse par, unde uterque quotiens par, et $2bb'+2\Delta,2bb'-2\Delta$ per $mm'$ divisibiles.

Iam ex undecim aequationibus fundamentalibus facile deducuntur sex sequentes:
\[
\begin{aligned}
APP&=aa'q'q'-2ab'qq'+ac'qq\\
AQQ&=aa'q''q''-2ba'qq''+ca'qq\\
ARR&=aa'q'''q'''-2(bb'+\Delta)qq'''+cc'qq\\
ASS&=ac'q''q''-2(bb'-\Delta)q'q''+ca'q'q'\\
ATT&=ac'q'''q'''-2bc'q'q'''+cc'q'q'\\
AUU&=ca'q'''q'''-2cb'q''q'''+cc'q''q''
\end{aligned}
\]

Hinc sequitur, omnes $APP,AQQ$ etc. divisibiles esse per $mm'$, unde facile derivatur, quoniam $kk$ divisor communis maximus numerorum $PP,QQ,$
